\begin{abstractpage}

	\begin{abstract}{romanian}
		Această lucrare analizează comparativ performanța unor metode clasice și a unor modele Deep Learning pentru segmentarea imaginilor la nivel de obiect.
		Sunt evaluate două metode clasice, de tip clustering, K-Means și Mean-Shift, precum și două rețele neuronale: U-Net și Slot Attention. Experimentele sunt realizate pe setul de date
		Tetrominoes din colecția Multi-Object Datasets, utilizând metrica Adjusted Rand Index (ARI) pentru evaluarea calității segmentării.

		Rezultatele experimentale arata faptul că metodele clasice, în special Mean-Shift Segmentation, obtin performanțe superioare modelelor de Deep Learning în contextul
		acestui set de date. În mod surprinzător, arhitectura Slot Attention, atât în varianta supervizată,
		cât și în cea nesupervizată, înregistrează cele mai scăzute scoruri, în ciuda resurselor de antrenare superioare. De asemenea, performanța modelului U-Net sugerează apariția fenomenului
		de supraînvățare, indicând o capacitate limitată de generalizare.

		Analiza rezultatelor sugerează că simplitatea și dimensiunea redusă a setului de date favorizează metodele clasice de segmentare, în timp ce modelele DeepLearning performează sub așteptari.
	\end{abstract}

	\begin{abstract}{english}
		This paper presents a comparative analysis of classical methods and Deep Learning models for object-level image segmentation.
		Two classical clustering methods, K-Means and Mean-Shift, as well as two neural network architectures, U-Net and Slot Attention,
		are evaluated. Experiments are conducted on the Tetrominoes dataset from the Multi-Object Datasets collection, using the
		Adjusted Rand Index (ARI) to assess segmentation quality.
		Experimental results indicate that classical methods, particularly Mean-Shift Segmentation, achieve superior performance compared to Deep Learning models on this dataset.
		Surprisingly, the Slot Attention architecture, both in its supervised and unsupervised variants, records the lowest scores despite having greater training resources.
		Additionally, the performance of the U-Net model suggests the presence of overfitting.
		The analysis suggests that the simplicity and small size of the dataset favor classical segmentation methods, while Deep Learning models underperform relative to expectations.
	\end{abstract}

\end{abstractpage}
